% WaRF Technical Reference
% Companion document to warf_talk.tex
% Build: pdflatex warf_technical.tex

\documentclass[11pt,a4paper]{article}

\usepackage{amsmath,amssymb}
\usepackage{booktabs}
\usepackage{geometry}
\usepackage{hyperref}
\usepackage{parskip}
\usepackage{microtype}
\usepackage{xcolor}
\usepackage{enumitem}

\geometry{margin=2.5cm}
\hypersetup{colorlinks=true,linkcolor=blue!60!black,citecolor=blue!60!black,urlcolor=blue!60!black}

\title{\textbf{WaRF: Weighted $n$-Agent Reliability Framework}\\[0.4em]
\large Technical Reference}
\author{Teymur Heydarov}
\date{August 2026}

\begin{document}
\maketitle
\tableofcontents
\bigskip\hrule
\bigskip

% =============================================================================
\section{Overview}

WaRF is a multi-agent code review framework that combines Bayesian weight
calibration with sequential hypothesis testing to produce a single
APPROVE/REJECT decision with a calibrated confidence level.

A pool of $n$ specialised LLM agents votes independently on the submitted
code. Each vote updates a cumulative evidence score $\Lambda$ via the
Sequential Probability Ratio Test (SPRT). The framework stops as soon as
$\Lambda$ crosses a decision boundary, or after all $n$ agents have voted.
Agent reliability weights $\omega \in [0,1]$ are per-agent and per-domain;
they are learned online from oracle-labelled reviews via a Bayesian update
rule.

% =============================================================================
\section{Agent Pool}

\subsection{Personas}

The default pool contains five LLM agents with fixed system prompts:

\begin{center}
\begin{tabular}{lll}
\toprule
Agent & Evaluation focus & Domain strengths \\
\midrule
AUDITOR      & Compliance audit, neutral         & Security, Logic \\
BOUNDARY     & Edge cases, missing guards        & Logic \\
ADVOCATE     & Author intent, false-positive resistance & General \\
SKEPTIC      & Adversarial reasoning, threat surface & Security \\
TEST-FOCUSED & Coverage, missing tests           & Logic \\
\bottomrule
\end{tabular}
\end{center}

Each agent receives: (1) a system prompt defining its persona, (2) the
code under review, and (3) an instruction to respond with exactly one word
(APPROVE or REJECT) on the first line, followed by free-text reasoning.
Agent responses are parsed for the vote; reasoning is logged and shown in
the dashboard but is not consumed by the stopping rule.

An optional sixth voter, \textbf{Bandit} (a deterministic SAST analyser),
is available in the SECURITY domain via the \texttt{--sast} flag.
It is assigned a fixed weight $\omega = 0.75$ and does not participate in
the Bayesian calibration cycle.

\subsection{Domain Classification}

Agents are indexed to three review domains:

\begin{description}[leftmargin=2.5cm,style=sameline]
  \item[\textsc{Security}] Authentication, injection, cryptography, input validation, serialisation.
  \item[\textsc{Logic}]    Correctness, edge cases, type safety, control flow, missing guards.
  \item[\textsc{Performance}] Algorithmic complexity, blocking calls, memory allocation, N+1 patterns.
\end{description}

The domain is specified by the caller via \texttt{--domain}.
Agent weights are domain-specific: an agent's LOGIC weight is entirely
independent of its SECURITY weight, capturing that BOUNDARY is more
reliable on logic defects than on security ones, and SKEPTIC is the
reverse.

% =============================================================================
\section{Sequential Probability Ratio Test}

\subsection{Cumulative Evidence Score}

After each agent vote, WaRF updates the log-likelihood ratio:

\begin{equation}
  \Lambda_k = \Lambda_{k-1} + s_k \cdot \omega_k \cdot \ell
  \label{eq:lambda}
\end{equation}

\noindent where:
\begin{itemize}
  \item $k$ is the vote index (agents are sampled in a fixed order)
  \item $s_k \in \{+1,-1\}$ is the vote sign (APPROVE $=+1$, REJECT $=-1$)
  \item $\omega_k$ is the calibrated weight of agent $k$ in the active domain
  \item $\ell \approx 1.253$ is the base log-likelihood ratio constant
  \item $\Lambda_0 = 0$
\end{itemize}

The LLR increment is asymmetric in the implementation. With $P_1 = 0.80$
and $P_0 = 0.30$ encoding prior beliefs about agent reliability:

\[
  \ell_{\text{APPROVE}} = \ln\!\frac{P_1}{P_0}
    = \ln\!\frac{0.80}{0.30} \approx +0.981
  \qquad
  \ell_{\text{REJECT}} = \ln\!\frac{1-P_1}{1-P_0}
    = \ln\!\frac{0.20}{0.70} \approx -1.253
\]

Because $|\ell_{\text{REJECT}}| > |\ell_{\text{APPROVE}}|$, unanimous REJECT
evidence accumulates faster than unanimous APPROVE evidence. In the remainder
of this document, $\ell \approx 1.253$ denotes the REJECT magnitude and is
used to reason about boundary reachability in the REJECT direction; the
APPROVE direction requires a larger pool or higher $\omega$ to compensate.

\subsection{Decision Boundaries}

WaRF uses asymmetric SPRT boundaries derived from target error rates:

\begin{align}
  A &= \ln\!\left(\frac{1-\beta}{\alpha}\right)
     = \ln\!\left(\frac{0.90}{0.05}\right) \approx +2.890
  & \text{(APPROVE boundary)} \label{eq:A} \\
  B &= \ln\!\left(\frac{\beta}{1-\alpha}\right)
     = \ln\!\left(\frac{0.10}{0.95}\right) \approx -2.251
  & \text{(REJECT boundary)}  \label{eq:B}
\end{align}

\noindent with $\alpha = 0.05$ (false-alarm rate) and $\beta = 0.10$
(miss rate). The asymmetry $|A| > |B|$ is deliberate: requiring stronger
evidence for APPROVE than for REJECT. A false rejection triggers human review,
which is recoverable. A false approval on a defective change is not.

\subsection{Stopping Rule}

After vote $k$:
\begin{itemize}
  \item If $\Lambda_k \geq A$: \textbf{stop}, decision $=$ APPROVE
  \item If $\Lambda_k \leq B$: \textbf{stop}, decision $=$ REJECT
  \item Otherwise: invite the next agent (or proceed to synthesis)
\end{itemize}

If the pool is exhausted without crossing a boundary, WaRF uses a synthesis
threshold $\theta = 0.2$:

\begin{equation}
  \text{decision} = \begin{cases}
    \textsc{Approve}      & \text{if } \Lambda_n > \theta \\
    \textsc{Reject}       & \text{if } \Lambda_n < -\theta \\
    \textsc{Inconclusive} & \text{if } |\Lambda_n| \leq \theta
  \end{cases}
\end{equation}

\subsection{Maximum Reachable Evidence and Pool Size}

With a pool of $n$ agents voting unanimously in one direction:

\begin{equation}
  |\Lambda|_{\max} = \ell \cdot \sum_{k=1}^{n} \omega_k
  \label{eq:max_lambda}
\end{equation}

SPRT early stopping is only possible if $|\Lambda|_{\max}$ exceeds the
relevant boundary. Table~\ref{tab:reachability} shows reachability for
representative configurations.

\begin{table}[h]
\centering
\begin{tabular}{lrrrll}
\toprule
Configuration & $n$ & $\sum\omega_i$ & $|\Lambda|_{\max}^{\text{rej}}$ & REJECT? & APPROVE? \\
\midrule
Uniform $\omega=0.5$           & 2 & 1.000 & 1.253 & No  & No \\
Top-2 current ($\omega{\approx}0.605$)  & 2 & 1.211 & 1.517 & No  & No \\
Top-3 current                   & 3 & 1.811 & 2.268 & \textbf{Yes} & No \\
Top-4 current                   & 4 & 2.332 & 2.922 & Yes & \textbf{Yes} \\
Full 5-agent pool               & 5 & 2.803 & 3.511 & Yes & Yes \\
\bottomrule
\end{tabular}
\caption{SPRT boundary reachability at current calibrated weights
($|\lambda_B|=2.251$, $\lambda_A=2.890$). $|\Lambda|_{\max}^{\text{rej}}$ uses
$\ell_{\text{REJECT}}=1.253$; the APPROVE direction (using $\ell_{\text{APPROVE}}=0.981$)
requires $n\geq 4$ at current weights.}
\label{tab:reachability}
\end{table}

At current calibrated weights, $n=3$ is the minimum pool that can reach the
REJECT boundary; $n=4$ for the APPROVE boundary. With $n=2$, the framework
always exhausts the pool and synthesizes, capped at \textsc{Medium} confidence.
These thresholds are weight-dependent: as agents accumulate more oracle labels
and $\hat\omega$ rises toward 1, smaller pools become sufficient.

% =============================================================================
\section{Bayesian Weight Calibration}

\subsection{Model}

Each agent $a$ maintains a per-domain Beta distribution over its reliability:

\begin{equation}
  \omega_a^{(d)} \sim \text{Beta}\!\left(\alpha_a^{(d)},\, \beta_a^{(d)}\right)
\end{equation}

The prior is $\alpha_0 = \beta_0 = 1$, the uniform distribution, corresponding
to $\omega = 0.5$ with maximum uncertainty (two pseudo-observations).

\subsection{Update Rule}

After an oracle-labelled review in domain $d$, for each agent $a$ that voted:
\begin{itemize}
  \item Vote matched oracle verdict: $\alpha_a^{(d)} \mathrel{+}= 1$
  \item Vote did not match oracle verdict: $\beta_a^{(d)} \mathrel{+}= 1$
\end{itemize}

The point estimate used in Equation~\ref{eq:lambda} is the posterior mean:

\begin{equation}
  \hat{\omega}_a^{(d)} = \frac{\alpha_a^{(d)}}{\alpha_a^{(d)} + \beta_a^{(d)}}
\end{equation}

Agents that are skipped by SPRT early stopping do not receive an update —
only agents that actually voted are evaluated against the oracle.

\medskip\noindent\textbf{Interpretation.} $\hat\omega$ is an estimate of
\emph{agent-domain fit}, not of the agent's absolute quality. An agent that
performs well on LOGIC code but poorly on SECURITY code will accumulate high
$\hat\omega^{\textsc{logic}}$ and low $\hat\omega^{\textsc{security}}$; neither
value characterises the agent in isolation. Changes in $\hat\omega$ over time
signal that the current corpus is shifting in ways that suit or penalise the
agent's evaluation heuristics — not that the agent itself has changed.

\subsection{Confidence Intervals}

A 95\% credible interval for $\omega$ is computed via the Normal
approximation to the Beta posterior:

\begin{equation}
  \hat{\omega} \;\pm\; 1.96\,\sqrt{\frac{\hat{\omega}(1-\hat{\omega})}
                                         {\alpha + \beta}}
\end{equation}

Wide intervals indicate few oracle-labelled reviews in the domain; the
weight is not yet reliable and should be interpreted with caution.

\subsection{Oracle Labelling}

The oracle is human judgment: after a review session, a developer marks the
ground-truth verdict. WaRF compares each agent's vote to the oracle and
updates weights accordingly. The system credits the correct \emph{outcome},
not the quality of reasoning: a correct vote for a spurious reason receives
the same reward as a correct vote for the right reason.

This is a deliberate black-box abstraction. Evaluating reasoning quality
would require a developer to read and judge every explanation — which is
exactly the work WaRF is meant to reduce.

% =============================================================================
\section{Confidence Classification}

WaRF assigns one of three confidence labels to each decision:

\begin{center}
\begin{tabular}{llll}
\toprule
Label & Condition & Empirical accuracy & Sessions \\
\midrule
\textsc{High}   & SPRT boundary crossed ($|\Lambda| \geq |\lambda_B|$ or $\geq \lambda_A$) & 75\%  & 12 \\
\textsc{Medium} & Pool exhausted, $|\Lambda_n| > \theta$                                    & 63\%  & 35 \\
\textsc{Low}    & Synthesis, $|\Lambda_n| \leq 2\theta$                                     & 27\%  & 15 \\
\bottomrule
\end{tabular}
\end{center}

\smallskip
Accuracy rises monotonically from LOW (27\%) to HIGH (75\%) across 62
oracle-labelled sessions, confirming that the SPRT boundary marks a
meaningful evidence threshold. The LOW anti-correlation (worse than random)
indicates that when the pool exhausts without accumulating much signal, the
synthesis decision is actively unreliable; such cases should always route to
human review.

The practical utility of the HIGH label is in \emph{selective automation}:
a downstream policy can auto-approve HIGH-confidence APPROVE decisions, block
on HIGH-confidence REJECT, and escalate MEDIUM or LOW to a human reviewer.
This stratification is only possible because the SPRT boundary produces a
qualitatively distinct evidence tier, not merely a higher vote count.

% =============================================================================
\section{Debate Round}

When a Round-1 decision has LOW or MEDIUM confidence, WaRF can optionally
invoke a debate round. Each agent receives all other agents' Round-1 votes
and reasoning and is asked whether it maintains or revises its position.

\subsection{Prompt Design}

The debate prompt is not neutral. Two framings produce systematically
different outcomes on the same code:

\begin{center}
\begin{tabular}{p{5.5cm}p{2.5cm}p{4.5cm}}
\toprule
Prompt & Typical outcome & Mechanism \\
\midrule
\textit{``Do you revise your position?''} (no anchor)
  & Conformity
  & Minority reads majority as evidence and capitulates \\[0.5em]
\textit{``Your Round-1 vote was X. Maintain or revise? Explain what others got wrong if you maintain.''}
  & Rational persuasion
  & Minority has a path to defend its position \\
\bottomrule
\end{tabular}
\end{center}

The anchoring prompt is the current default in WaRF. The conformity nudge
was an early design that was found to amplify the wrong majority — including
with HIGH confidence after SPRT fired on the conforming votes.

\subsection{When Debate Helps and When It Does Not}

Debate propagates correct reasoning that agents can evaluate and find
convincing. It does not cure false beliefs held by an agent whose training
does not include the relevant ground truth.

\begin{center}
\begin{tabular}{p{3.5cm}p{2.0cm}p{5.5cm}}
\toprule
Case & Outcome & Reason \\
\midrule
\texttt{sessions\_pre\_auth\_strip.py}
  & \textcolor{green!60!black}{Corrected}
  & Correct minority (SKEPTIC) had a verifiable argument about \texttt{rebuild\_auth} stripping; majority evaluated and updated \\[0.5em]
\texttt{hooks.py}
  & No change
  & Incorrect agent (SKEPTIC) held a false belief about \texttt{collections.abc.Callable}; correct explanation did not penetrate \\[0.5em]
\texttt{\_internal\_utils.py}
  & \textcolor{red}{Worsened}
  & Plausible wrong argument (``no None guard'') captured the correct minority; SPRT fired on the wrong unanimous verdict \\
\bottomrule
\end{tabular}
\end{center}

\textbf{Key principle: persuasion $\neq$ correctness.} Debate propagates
the most \emph{persuasive} argument, not the correct one. A structurally
sound generic heuristic (``no None guard'') can capture the correct minority
even when it is inapplicable in context. When this happens, the result is a
unanimous wrong verdict with HIGH confidence — the worst failure mode, because
the SPRT boundary provides false assurance. The oracle is the only in-system
corrective; debate has no mechanism to distinguish a persuasive argument from
a true one.

% =============================================================================
\section{Empirical Results}

Results are drawn from 246 sessions across PSF/requests library code and
manually constructed validation files (M-00 through M-13).

\subsection{Calibration}

On 62 oracle-labelled sessions (reliability diagram):

\begin{center}
\begin{tabular}{lccc}
\toprule
Confidence & Sessions & Correct & Accuracy \\
\midrule
\textsc{High}   & 12 &  9 & 75.0\% \\
\textsc{Medium} & 35 & 22 & 62.9\% \\
\textsc{Low}    & 15 &  4 & 26.7\% \\
\midrule
Overall         & 62 & 35 & 56.5\% \\
\bottomrule
\end{tabular}
\end{center}

The 48-percentage-point spread from LOW (27\%) to HIGH (75\%) validates
that the SPRT boundary marks a meaningful evidence threshold. Overall
accuracy of 56.5\% reflects a corpus balanced between straightforward and
hard cases; the confidence-stratified numbers are the informative signal.
Compared to the previous corpus snapshot, the MEDIUM band has grown (22
to 35 sessions) while LOW has shrunk (28 to 15), consistent with agents
accumulating enough calibration history to escape the inconclusive range
on more reviews.

\subsection{Token Efficiency}

SPRT early stopping fires on approximately 44\% of sessions (109 of 246).
On those sessions the average token usage drops substantially compared to a
full-pool run:

\begin{center}
\begin{tabular}{lrll}
\toprule
Configuration & Tokens & Confidence & Notes \\
\midrule
5-agent, pool exhausted        & $\sim$16\,100 & \textsc{Medium} (max) & All 5 agents called \\
5-agent, SPRT at $k=3$         & $\sim$3\,200  & \textsc{High}         & 2 agents skipped \\
2-agent, any outcome           & $\sim$4\,000  & \textsc{Medium} (max) & Boundaries unreachable \\
\bottomrule
\end{tabular}
\end{center}

Token cost is more strongly driven by file size than pool size; a 5-agent
run on a small file is cheaper than a 2-agent run on a large one. The early
stopping rate of 44\% has grown with corpus size: as agents accumulate
oracle labels and $\hat\omega$ converges away from the uninformative prior,
evidence accumulates faster per agent, making boundary crossings more frequent.

\subsection{2-Agent Pool Comparison}

The top-2 agents by current calibrated $\hat\omega$ (ADVOCATE and AUDITOR,
both $\approx 0.605$) are structurally capped at \textsc{Medium} confidence
(Section~3.4, Table~1): their combined $|\Lambda|_{\max} \approx 1.52$, well
below the REJECT boundary at 2.251. This cap holds for any choice of two
agents at current weight distributions.

The critical practical consequence is \emph{selective automation}. A downstream
CI policy can only stratify decisions if confidence levels are qualitatively
distinct: HIGH to auto-approve or auto-block; MEDIUM and LOW to escalate.
With a 2-agent pool, every decision returns MEDIUM regardless of how
unambiguous the evidence is. The routing signal disappears. The 5-agent pool
produces HIGH on 44\% of reviews — those are the sessions that can be handled
without human intervention.

Counter-argument to \cite{single2026}: their finding concerns mathematical
reasoning tasks, where the defect is derived from premises. Code review is
different: the defect is present or absent in the artefact, and the ensemble's
value is not deliberation but accumulating the evidence mass required to cross
the SPRT boundary with HIGH confidence. The 2-agent pool achieves comparable
accuracy on average but never produces the HIGH-confidence tier that makes
selective automation possible.

% =============================================================================
\section{Known Limitations}

\begin{description}
\item[Type 1 — Persona bias.]
  Agents with affirmative personas (ADVOCATE, AUDITOR) evaluate whether the
  code does what its author intended — not whether what the author intended is
  sufficient. They correctly assess that the logic is coherent and then stop.
  They do not ask what happens when an integer overflows, or when a pointer is
  null, because those scenarios lie outside the framing their persona uses to
  read code. BOUNDARY and TEST-FOCUSED partially correct this by design, but
  any pool dominated by affirmative personas inherits the bias structurally.
  More agents of the same persona type amplify rather than correct it.

\item[Type 2 — Shared knowledge gap.]
  When all agents share the same training-level blind spot, WaRF produces
  a unanimous wrong verdict with HIGH confidence. CVE-2024-47081 is the
  canonical example: all five agents approved the vulnerable \texttt{urlparse}
  hostname parsing line. Adding more agents with the same training coverage
  amplifies the wrong signal. Bandit (SAST) partially mitigates this for
  patterns in its ruleset.

\item[Type 3 — False-consensus capitulation.]
  In debate, a correct minority can be captured by a plausible wrong majority
  argument (see \texttt{\_internal\_utils.py} case). Once SPRT fires on the
  conforming votes, there is no in-session recovery; the oracle update is the
  only corrective mechanism.

\item[Boundary reachability.]
  At current calibrated weights, $n \geq 3$ is required to reach the REJECT
  boundary and $n \geq 4$ for the APPROVE boundary (Section~3.4, Table~1).
  With $n=2$, the framework always exhausts the pool and synthesizes at
  MEDIUM. As agents accumulate more oracle labels and $\hat\omega$ rises,
  smaller pools may become sufficient; conversely, a domain with few labelled
  sessions keeps $\hat\omega$ near the prior (0.5), making boundaries harder
  to reach. Adding more agents increases cost proportionally whenever the
  boundary is not crossed before pool exhaustion.

\item[Oracle dependency.]
  Bayesian calibration requires oracle-labelled reviews. Without labelling,
  all weights remain at their prior ($\hat\omega = 0.5$) and the SPRT
  reduces to an unweighted vote counter. The quality of calibration is bounded
  by the labelling rate.
\end{description}

% =============================================================================
\section{Related Work}

\textbf{ConSol}~\cite{consol2025} applies SPRT to a single language model at
varied sampling temperature to determine when the model has converged on a
consistent answer. WaRF's extension: heterogeneous agent personas with
domain-specific learned weights. Temperature variance gives noise around one
perspective; different personas give genuinely different evaluation heuristics
(compliance audit vs.\ adversarial probing). The domain-specific $\omega$
weights are only meaningful because the agents are actually different.

\textbf{Asymptotically Optimal Sequential Testing}~\cite{optimal2026} derives
optimal stopping policies for LLM ensembles under a combined cost-plus-wait
objective. The optimal policy uses at most two LLMs. The objective differs
from WaRF's: WaRF optimises calibrated accuracy at minimum agent calls without
a cost term, and does not assume homogeneous agents.

\textbf{Wisdom and Delusion of LLM Ensembles}~\cite{wisdom2025} demonstrates
that ensembles do not always outperform the single best agent, particularly
when agents share training data and therefore share blind spots. This is
consistent with WaRF's Type 2 failure mode and validates the distinction
between \emph{persona diversity} (which WaRF provides) and
\emph{knowledge diversity} (which requires external tools or data augmentation).

\textbf{Single-Agent Outperforms Multi-Agent Under Equal Token Budget}~\cite{single2026}
finds that on mathematical reasoning tasks, adding agents under a fixed budget
degrades performance. Code review differs: the defect is not derived from
premises but observed in the artefact, and the ensemble's value lies in
accumulating sufficient evidence mass to cross the SPRT boundary — not in
deliberation.

\begin{thebibliography}{9}
\bibitem{consol2025}
  ConSol: Sequential Consistency for LLM Outputs via SPRT.
  arXiv:2503.17587, 2025.
\bibitem{optimal2026}
  Asymptotically Optimal Sequential Testing for LLM Ensembles.
  arXiv:2604.01086, 2026.
\bibitem{wisdom2025}
  The Wisdom and Delusion of LLM Ensembles.
  arXiv:2510.21513, 2025.
\bibitem{single2026}
  Single-Agent Outperforms Multi-Agent Under Equal Token Budget.
  arXiv:2604.02460, 2026.
\end{thebibliography}

\end{document}
